\appendix

\section{The constants and the numerical steps of the shape bound}\label{app:shape}

\subsection{The constants of the flag bound}

For $n=15$ we have $B_r=\bigl(r/(8(16-r))\bigr)^2$. Table~\ref{tab:cr} lists
$c_r=B_r-B_{r-1}$ and the ratios $c_{r+1}/c_r$. All entries are exact. The
product of the fourteen values $c_r$ is the constant $K_{15}$ of
Section~\ref{sec:shape}, and
\begin{align*}
 15\log\frac{A_{15}}{30}&=-\frac{15}2\log K_{15}-15\log30\\
 &=497.0298125704945\ldots ,
\end{align*}
computed with fifty significant digits. Since
\[
 \log720720=13.4880059913253\ldots ,
\]
the logarithms of the primes $p\ge17$ must add up to more than
\[
 483.5418065791692\ldots .
\]

\begin{table}[t]
\caption{The increments $c_r=B_r-B_{r-1}$ of the prefix bounds for fifteen
speeds and their ratios. The least ratio is at $r=6$.}
\label{tab:cr}
\centering
\small
\begin{tabular}{@{}rlll@{}}
\toprule
$r$ & $c_r$ & $c_{r+1}/c_r$ & \\
\midrule
\CRROWS
\bottomrule
\end{tabular}
\end{table}

\subsection{The inequalities in the proof of Lemma~\ref{lem:shape}}

Every inequality below compares two rational numbers, after raising both
sides to a power that clears the fractional exponents; the decimal values
are given to show the margins.

\emph{The root $\alpha$.} The polynomial $g(u)=14u^2-13u+1$ takes the
values
\begin{gather*}
 g(0.0846)=0.00040024>0,\\
 g(0.0847)=-0.00066274<0,
\end{gather*}
so the smaller root $\alpha$ of $g$ lies in $(0.0846,0.0847)$.

\emph{The sign of $P$.} For $0<u<\alpha<0.0847$ the terms $7u^5+14u^4$ are
positive and $-4u-11u^2-27u^3$ is decreasing, so
\begin{align*}
 P(u)&>1-4(0.0847)-11(0.0847)^2-27(0.0847)^3\\
 &=0.565878\ldots>0 .
\end{align*}

\emph{The lower bound for $a(\alpha)$.} Since $q$ increases on $(0,1)$ and
$u^{-2/15}$ decreases, $a(\alpha)>q(0.0846)/0.0847^{2/15}$. Raising to the
fifteenth power, the claim $q(0.0846)/0.0847^{2/15}>\frac85$ becomes
\begin{align*}
 q(0.0846)^{15}&=9.51331\ldots\\
 &>\Bigl(\frac85\Bigr)^{15}(0.0847)^2=8.27116\ldots .
\end{align*}

\emph{Two or more coordinates equal to one.} The ratio of consecutive
values of $f(k)=\frac k2\,2^k\bigl(\frac85\bigr)^{15-k}$ is
$f(k+1)/f(k)=\frac{k+1}k\cdot\frac54>1$, so $f$ increases, and
$f(2)=4\bigl(\frac85\bigr)^{13}=1801.43\ldots$.

\emph{The root $t_0$.} The cubic $378t^3+25t^2+10t-1$ takes the values
$-0.0015037\ldots$ at $t=0.0724$ and $0.00045378\ldots$ at $t=0.0725$.

\emph{The value $H(t_0)$.} The factors $1+2t+27t^2$ and $q(t)$ increase
and $t^{-28/15}$ decreases, so on $[0.0724,0.0725]$
\[
 H(t)\ \ge\ \frac{(1+2\cdot0.0724+27\cdot0.0724^2)\,q(0.0724)^{13}}
 {0.0725^{28/15}} .
\]
Raising to the fifteenth power, the right side exceeds $942.5$ because
\begin{align*}
 (1+2\cdot0.0724&+27\cdot0.0724^2)^{15}\,q(0.0724)^{195}\\
 &=5.05703\ldots\cdot10^{12},\\
 942.5^{15}\cdot0.0725^{28}&=5.05457\ldots\cdot10^{12}.
\end{align*}
The bound itself is $942.53\ldots$, against the true minimum
$H(t_0)=944.96\ldots$ at $t_0=0.0724768\ldots$.

\emph{The extremal shape.} The minimum of $H$ is attained at the point
$u=(1,t_0,\dots,t_0)$, where
\[
 R_{15}=\bigl(225/H(t_0)\bigr)^{1/2}=0.48796\ldots .
\]
Using $\frac12$ in place of this value in Theorem~\ref{thm:flag} weakens
the logarithmic bound by $15\log(0.5/0.48796)=0.366$.

\subsection{The algebraic identities}

The identities used in the proof are checked by expanding both sides as
rational functions: the formula~\eqref{eq:phiprime} for $\varphi'$, the
identity $14q(u)+g(u)=15(1+u)$, the formula for $L'$, the logarithmic
derivative of $a$, the closed form of $H(t)$ on the line
$u=(1,t,\dots,t)$ and its logarithmic derivative. For the last one,
\[
 \frac{d}{dt}\log H(t)=\frac{2+54t}{1+2t+27t^2}+\frac{13(2-2t)}{q(t)}
 -\frac{28}{15t},
\]
and putting the three terms over the common denominator
$15t\,q(t)(1+2t+27t^2)$ gives the numerator
$-28(t-1)(378t^3+25t^2+10t-1)$; since $q(t)=-(t^2-2t-1)$, this is the
expression in the proof.
